% !TeX root = ../../infdesc.tex
\section{Limits of functions}
\secbegin{secLimitsOfFunctions}


\begin{exercise}
Let $a,b \in \mathbb{R}$ with $a<b$. Prove that $\overline{(a,b)} = \overline{(a,b]} = \overline{[a,b)} = [a,b]$.
\end{exercise}

\begin{exercise}
Let $f : D \to \mathbb{R}$ be a function, let $a \in D$ and let $\ell \in \mathbb{R}$. Fix some sequence $(x_n)$ of elements of $D$, not eventually constant, such that $(x_n) \to a$. Prove that if $f(x) \to \ell$ as $x \to a$, then the sequence $(f(x_n))$ converges to $\ell$.
\end{exercise}


\begin{exercise}
\label{exLimitsOfFunctionsAreUnique}
Let $f : D \to \mathbb{R}$ be a function, let $a \in D$, and $\ell_1, \ell_2 \in \mathbb{R}$. Prove that if $f(x) \to \ell_1$ as $x \to a$, and $f(x) \to \ell_2$ as $x \to a$, then $\ell_1=\ell_2$.
\end{exercise}


\begin{exercise}
\label{exInfiniteLimitsOfFunctionsAreUnique}
Let $f : D \to \mathbb{R}$ be a function and $\ell_1, \ell_2 \in \mathbb{R}$. Prove that if $D$ is unbounded above, and if $f(x) \to \ell_1$ as $x \to \infty$ and $f(x) \to \ell_2$ as $x \to \infty$, then $\ell_1=\ell_2$. Prove the analogous result for limits as $x \to -\infty$ in the case when $D$ is unbounded below.
\end{exercise}

